% !TeX root = ../Book1.tex
\begin{figure}[htbp]
	\centering
	\begin{tikzpicture}[scale=0.95,font=\small]
		\begin{scope}[xshift=-3cm]
			\draw[gray,dashed] (0,0) circle (1.6);
			\draw[blue!70!black,thick,->] (0,0)--(1.05,0);
			\draw[red!70!black,thick,->] (0,0) .. controls (-0.1,1.05) and (1.4,1.2) .. (1.05,0);
			\fill (0,0) circle (1.6pt) node[below left] {\(p\)};
			\fill (1.05,0) circle (1.6pt) node[below right] {\(q\)};
			\node at (0,-2) {留在 \(V\) 中的曲线};
		\end{scope}
		\begin{scope}[xshift=3cm]
			\draw[gray,dashed] (0,0) circle (1.6);
			\draw[blue!70!black,thick,->] (0,0)--(1.05,0);
			\draw[red!70!black,thick,->] (0,0) .. controls (-0.4,1.4) and (0,2.3) .. (1.4,2) .. controls (2.7,1.7) and (2.1,-0.15) .. (1.05,0);
			\fill (0,0) circle (1.6pt) node[below left] {\(p\)};
			\fill (1.05,0) circle (1.6pt) node[below right] {\(q\)};
			\node at (0,-2) {先到达 \(r=\rho\) 的曲线};
		\end{scope}
	\end{tikzpicture}
	\caption[正常邻域中的最短路比较]{测地坐标中的两类竞争曲线。蓝线的长度为 \(r(q)\)；左图中角向运动会增加长度，右图的曲线仅走到虚线边界就至少需要长度 \(\rho\)。图中纸面上的一般距离不代表曲面距离。}
	\label{g1:fig:normal-neighborhood}
\end{figure}
